\subsection{3. Noch wichtiger: Ein Überlappungsmatrixelement ist kein
Transport}\label{noch-wichtiger-ein-uxfcberlappungsmatrixelement-ist-kein-transport}

Für dieselben zwei Moden seien

\codeword{|A>=a†|0>}, \codeword{|B>=b†|0>}.

Dann ist ihre Gram-Matrix

\codeword{S=[[1,1/√2],[1/√2,1]]}.

Wähle den vollkommen trivialen Hamiltonoperator \codeword{H₀=ωN}, mit
\codeword{N=c₁†c₁+c₂†c₂}. Die Matrix der Hamilton-Matrixelemente in
diesen beiden Chartvektoren ist

\codeword{K\_ij=<i|H₀|j>=ω S\_ij}.

K hat also eine nichtverschwindende Offdiagonale, obwohl auf dem
gesamten Einteilchenraum nur dieselbe Phase
\codeword{e\textasciicircum{}(−iωt)} entsteht. Das korrekte Eigenproblem
ist \codeword{Kv=E Sv}, nicht \codeword{Kv=Ev}. Es liefert
ausschließlich die entartete Energie ω. Die Wahrscheinlichkeit

\codeword{|<B|exp(−itH₀)|A>|²=1/2}

ist für jede Zeit gleich groß. Es hat keine Übertragung stattgefunden;
die Hälfte war schon bei t=0 als statische Überlappung vorhanden.

Das ist unmittelbar für den bisherigen nativen TFPT-Pol relevant. Dieser
besitzt im erklärten N=63-Vertrag eine einzelne Energie E\_h auf einem
64-dimensionalen Polraum:

\codeword{P\_h H P\_h = E\_h P\_h}.

Sind A und B lediglich zwei Beschreibungen oder Sonden innerhalb
desselben Polraums, gilt wiederum \codeword{K=E\_h S}: ein gemeinsamer
Phasenfaktor, keine durch H verursachte Ausbreitung zwischen den Marken.
Diese Aussage braucht keine große Diagonalisierung. Sie folgt allein aus
der bereits bewiesenen Entartung.

Ein passiver Chartwechsel verändert nur die Beschreibung. Ein aktiv
implementierter Rotationsoperator kann Zustände verändern; dann müssen
jedoch genau dieser Operator, seine physische Verfügbarkeit, Zeit- oder
Ressourcenskala und sein Instrument aus der Quelle nachgewiesen werden.
Das ist nicht durch den Namen »Intertwiner« erledigt.

Die mögliche globale Überlappungskonstruktion ist dadurch nicht
ausgeschlossen. Sie müsste tatsächliche zusätzliche globale Dynamik
beziehungsweise eine Bandaufspaltung aus der Quelle zeigen und die
lokalen Polbeweise darin erneut absichern.

\subsection{4. Chartwechsel, Verbindungen und Krümmung sind verschiedene
Daten}\label{chartwechsel-verbindungen-und-kruxfcmmung-sind-verschiedene-daten}

Wenn mehrere Charts lediglich vollständige orthonormale Rahmen R\_x
desselben festen Vektorraums sind, lauten die Basiswechsel

\codeword{U\_xy=R\_x† R\_y}.

Dann teleskopiert jedes Dreiecksprodukt:

\codeword{U\_12 U\_23 U\_31=I}.

Das gilt auch bei nichtkommutierenden R\_x; der Checker prüft ein
konkretes solches Beispiel. Durch bloße lokale Basiswahl entsteht daher
nicht automatisch ein physisch gekrümmtes Eichfeld. Auf einem Bündel
beschreiben Übergangsfunktionen seine Verklebung; eine Verbindung
enthält zusätzliche Paralleltransportdaten. Eine nichttriviale
Bündeltopologie ist ebenfalls nicht dasselbe wie bereits gewählte lokale
Krümmung.

Bei tatsächlich variierenden \textbf{echten Unterräumen} ist die
Situation interessanter. Überlappungen \codeword{ι\_x†ι\_y} müssen dann
nicht unitär sein. Drei Strahlen \codeword{(1,0)}, \codeword{(1,1)/√2},
\codeword{(1,i)/√2} besitzen das nichtreelle Schleifenprodukt
\codeword{(1+i)/4}. Das liefert eine geometrische Phase, aber sein
Betragsquadrat ist 1/8 und nicht eins. Eine solche Unterraumgeometrie
kann Ansatzpunkt für eine geometrische Verbindung sein; sie ist kein
automatisch ausgeführter normerhaltender Transport und keine schon
gewonnene Eichfeldkinetik.

Zudem legt eine interne Eichverbindung alleine keine Raumzeitmetrik
fest. Verschiedene Linkphasen können auf demselben Graphen bei denselben
Operationskosten existieren. »Die U\_xy ändern sich« bedeutet ohne
weiteren Nachweis nicht bereits »die Raumzeitgeometrie ändert sich«.

\subsection{5. Zeit: Ein Grundzustand bewegt sich unter seinem H nicht
beobachtbar}\label{zeit-ein-grundzustand-bewegt-sich-unter-seinem-h-nicht-beobachtbar}

Die im Anhang skizzierte Entwicklung \codeword{Ω -> exp(−itH)Ω} erzeugt
bei einem Energieeigenzustand lediglich \codeword{exp(−itE₀)Ω}. Der
Dichteoperator bleibt exakt gleich. Der Checker bestätigt das an einem
kleinen Fockzustand.

Damit wird weder die Zeitordnung noch ein Zeitpfeil erzeugt.
Nichtstationäre Zustände und mehrzeitige Korrelationsfunktionen können
natürlich eine Dynamik anzeigen; eine gerichtete Record- oder
Präparationsgeschichte braucht ihre eigenen Bedingungen. Auch eine
relationale Zeit aus bedingten Zuständen wäre als Forschungsansatz
möglich, verlangte aber eine explizite Uhr, ihren gemeinsamen Zustand
mit dem System und eine Konditionierungsregel. Der bloße globale
Eigenzustandsphasenfaktor reicht dafür nicht.

Die nun geprüfte endliche Clock kann eine aktive Operation auf inneren
Marken darstellen, falls sie physisch gewährt ist; dass sie zur
Quellsymmetrie gehört, macht sie nicht automatisch identisch mit der
Hamiltonzeit.

\subsection{6. Der binäre Index muss als Freiheitsgrad bewiesen
werden}\label{der-binuxe4re-index-muss-als-freiheitsgrad-bewiesen-werden}

Die skalare Hilfsdublett-Reparatur braucht zwei unabhängige,
gleichgeladene Komponenten, auf denen eine nichtentartete
antisymmetrische Form wirkt. Eine Orientierung der Überlappungsgeometrie
\textbf{könnte} so etwas liefern, tut es aber nicht allein durch das
Vorhandensein zweier Bezeichnungen.

Die drei einfachen Verwechslungen sind:

\begin{itemize}
\tightlist
\item
  Zwei Namen oder ±-Vorzeichen für dieselbe Mode bilden nur eine
  eindimensionale, redundante Beschreibung. Der Rückzug der
  antisymmetrischen Zweiform auf diesen Raum ist null.
\item
  Ist Rückwärts die Adjungierte der Vorwärtsoperation, haben f und f†
  entgegengesetzte U(1)-Ladungen. Das ist kein gleichgeladenes
  Hilfsdublett. Ihr bilinearer Kanal ist neutral, nicht von Ladung −2.
\item
  Bedeutet Vorwärts/Rückwärts links-/rechtshändige Weylfelder, wurden
  zwei verschiedene Lorentzdarstellungen eingeführt. Das ist nicht die
  zusätzliche gleichartige Hilfskopie des bereits geprüften
  Tensorvertrags.
\end{itemize}

Eine geometrisch hergeleitete Zweifachheit könnte also eine gute
Erklärung des Hilfsindex sein. Nötig wären zwei tatsächliche unabhängige
Kanäle, ihre CAR-/Ladungsrelationen, die Spin-Lorentz-Wirkung und die
nichtverschwindende Kopplung auf demselben Quellraum. Ein
Orientierungsbit liefert weder automatisch Weylspin noch die
Transformationsregel unter einer vollen Lorentzdrehung.

\subsection{7. Operationsgeometrie benötigt mehr als
Erreichbarkeit}\label{operationsgeometrie-benuxf6tigt-mehr-als-erreichbarkeit}

Minimale Kosten erfüllen unter geeigneter Komposition eine
Dreiecksungleichung. Ohne reversible Operationen mit symmetrischen
Kosten ist die resultierende Funktion jedoch nur eine gerichtete
Distanz; der Checker liefert den gerichteten Dreierzyklus mit
\codeword{d(A,B)=1}, \codeword{d(B,A)=2}. Unerreichbarkeit ergibt
unendliche Distanzen. Kostenfreie Chartwechsel können verschiedene
Beschreibungen auf Distanz null setzen; dann muss zunächst nach
physischer Äquivalenz quotiert werden.

Wachstum \codeword{V(R)\textasciitilde{}R³} ist ein Nachweis einer
dreidimensionalen \textbf{Wachstumsdimension im gewählten Kostenmaß},
nicht automatisch einer glatten dreidimensionalen Mannigfaltigkeit,
einer 3+1D-Lorentzmetrik oder einer universellen Lichtgeschwindigkeit.
Schon ein kubisches Gitter mit Laplace-Hamiltonoperator hat dieses
Volumenwachstum, aber bei kleinen Impulsen

\codeword{E(k)=Σ\_j(2−2cos k\_j) \textasciitilde{} |k|²},

nicht eine lineare relativistische Dispersion. Der eindimensionale
Taylor-Koeffizient dieser separierbaren Gegenkonstruktion wird exakt
geprüft.

Eine lokale Spektrallücke ist ebenso noch keine relativistische
Teilchenmasse. Dazu braucht es die gemeinsame Energie-/Impulsdeutung,
den entsprechenden Dispersionszweig und den Ladungs-/Referenzvertrag.
»Kopieren« sollte bei Teilchenpropagation nur bildlich verwendet werden:
kohärenter Transfer erzeugt nicht zwei unabhängige Kopien eines
unbekannten Quantenzustands.

\subsection{8. Gravitation bleibt eine dynamische
Verpflichtung}\label{gravitation-bleibt-eine-dynamische-verpflichtung}

Die linearen Fluktuationen einer aus der Quelle bestimmten globalen
Geometrie zu prüfen, ist ein sinnvoller Forschungsauftrag. Eine
transversale spurfreie Tensorzerlegung allein garantiert aber weder
einen masselosen physikalischen Spin-2-Pol noch positive Norm, genau
zwei Helizitäten oder universelle Kopplung. Selbst ein gesunder linearer
Spin-2-Sektor ersetzt nicht den Nachweis seiner nichtlinearen
Eich-/Zwangsstruktur und konsistenten Materiekopplung.

Das vorgeschlagene wechselseitige Schema Materie -\textgreater{}
bevorzugte Überlappungen -\textgreater{} Geometrie ist daher ein
mögliches Wirkungsprinzip, noch keine aus dem vorhandenen nativen H
gewonnene Gleichung.

\subsection{9. P versus NP: eine konkrete
Korrektur}\label{p-versus-np-eine-konkrete-korrektur}

Der Anhang sagt, P≠NP würde einen intrinsisch exponentiellen Suchaufwand
bedeuten. Das ist falsch. P≠NP würde ausschließen, dass \textbf{jedes}
NP-Entscheidungsproblem deterministisch in Polynomialzeit gelöst werden
kann. Es folgt daraus keine exponentielle untere Schranke: Eine
hypothetische Laufzeit wie \codeword{2\textasciicircum{}(√n)} ist
superpolynomiell und zugleich subexponentiell. Stärkere
Exponentialzeitaussagen benötigen zusätzliche Sätze beziehungsweise
Hypothesen.

NP bezieht sich zudem auf polynomial lange, polynomial prüfbare
Zertifikate in der Länge einer wohldefinierten Eingabekodierung.
Beliebige Erreichbarkeit in einem knapp beschriebenen Graphen mit
möglicherweise exponentiell langen Wegen ist nicht allein deshalb ein
NP-Problem. Eine geometrische Umformulierung muss Eingabelänge,
Zertifikatlänge, erlaubte Operationen und deren Kosten erhalten. Eine
Pfadmetapher löst diese Verpflichtungen nicht.

\subsection{10. Ein korrigierter, wirklich entscheidbarer
Anschluss-Test}\label{ein-korrigierter-wirklich-entscheidbarer-anschluss-test}

Ein sinnvoller nächster Versuch besteht nicht nur aus einem
unbeschrifteten \codeword{P U\_AB P}. Er sollte zusammen liefern:

\begin{enumerate}
\def\labelenumi{\arabic{enumi}.}
\tightlist
\item
  \textbf{Globale Quelle und Einbettungen:} konkrete primitive Algebra
  und zwei CAR-/Tensor-erhaltende Abbildungen; alle Kreuzrelationen und
  die tatsächliche globale Graduierung.
\item
  \textbf{Gemeinsamer Zustand:} derselbe globale H und Zustand;
  Nachweis, welche bisherigen lokalen Pole und Lücken darin noch gelten.
  Kein stiller Rückgriff auf den Produktgrundzustand unabhängiger
  Banken.
\item
  \textbf{Physische Operation:} aus einem erklärten Quellenwort
  erzeugter aktiver Operator oder Generator, mit Zeit- beziehungsweise
  Ressourcenskala. Passive Chartwechsel bleiben getrennt.
\item
  \textbf{Gram-bereinigte Dynamik:} \codeword{S\_ij=<h\_i|h\_j>} und
  \codeword{K\_ij=<h\_i|H|h\_j>} berechnen; das generalisierte
  Eigenproblem beziehungsweise eine orthonormalisierte Darstellung
  verwenden. Prüfen, ob mehr als \codeword{K=E\_h S} entsteht.
\item
  \textbf{Operativer Transfer:} dieselbe Anfangspräparation und
  Zielmessung; zeitabhängige Änderung gegenüber der statischen
  Überlappung und Fehlerkontrolle gegenüber dem übrigen Zustandsraum.
  Erst dann mit dem bekannten bedingten Zwei-Banken-Transfer
  vergleichen.
\end{enumerate}

Das nimmt die mögliche gemeinsame Herkunft ernst, korrigiert aber die
falschen Automatismen. Der stärkste sofortige Erkenntnisgewinn dieses
Audits ist negativ und präzise: \textbf{Allein durch Überlappung oder
Umbenennung eines entarteten nativen Polraums entsteht keine
Transportdynamik.} Ob der Compiler darüber hinaus genau die passende
globale Dynamik besitzt, ist die konkrete offene Frage.

\begin{center}\rule{0.5\linewidth}{0.5pt}\end{center}

\section{Anhang E: arithmetischer
Schleifenvorschlag}\label{anhang-e-arithmetischer-schleifenvorschlag}

\section{Zusatzprüfung: Schleifen, Primzahlen und der
Universalraum-Vorschlag}\label{zusatzpruxfcfung-schleifen-primzahlen-und-der-universalraum-vorschlag}

\begin{enumerate}
\def\labelenumi{\arabic{enumi}.}
\setcounter{enumi}{14}
\tightlist
\item
  September 2026, v1.6.6. Quellenprüfung des nachgereichten Textes
  \codeword{59dc0059-8914-48ca-953d-85933f66e00b}, kein neuer RH-Beweis
  und kein neuer arithmetischer Quellenoperator. Seine Aussagen über
  gemeinsame lokale Beschreibungen werden getrennt im Überlappungsaudit
  behandelt.
\end{enumerate}

\subsection{Urteil}\label{urteil-1}

\textbf{Ein globaler Wegraum ist eine sinnvolle Suchklasse; ein
Eulerprodukt über primitive Wege ist noch nicht das Eulerprodukt der
Riemannschen Zetafunktion.} Die im Text vorgeschlagene Xi-Determinante
bleibt ein hinreichendes Ziel unter expliziten Operatorannahmen. Das
Umbenennen ihres noch fehlenden Operators in „Generator der primitiven
Schleifen`` konstruiert ihn nicht.

\subsection{1. Was die endliche Bank tatsächlich
ausschließt}\label{was-die-endliche-bank-tatsuxe4chlich-ausschlieuxdft}

Der beibehaltene Beweis aus v1.6.5 gilt für den ursprünglichen
Hamiltonoperator mit 64 Fermion- und 60 Bosonmoden, Δ\textgreater0. Der
Hilbertraum ist wegen der Bosonen \textbf{nicht endlichdimensional}.
Endlich ist die Zahl der Moden. Mit \codeword{C=960g²/Δ} folgt aus
quadratischem Ergänzen

\[\adjustbox{max width=\linewidth}{$\displaystyle 
H\ge\frac\Delta2N_b-C,
\qquad
N_H(E)\le2^{64}\binom{\lfloor2(E+C)/\Delta\rfloor+60}{60}.
$}\]

Dies widerspricht einem vollständigen energieerhaltenden Spektrum
\codeword{E\_* log n + E\_off}, E\_*\textgreater0, weil dessen
Zählfunktion exponentiell wächst. Es ist kein Ausschluss beliebiger
arithmetischer Untersektoren oder eines anderen Operators mit
nichtlinearer Energiezuordnung.

Ein großer Wegraum kann eine andere Zählfunktion besitzen. Die
Gleichsetzung „mehr Wege = mehr orthogonale Zustände unter derselben
Energie`` muss aber bewiesen werden. Verschiedene Wörter können
denselben Operator oder Zustand darstellen; ein Quantenüberlagerungsraum
ist nicht ohne Weiteres der freie Wortraum. Bei unendlich vielen
identischen Zellen kann stattdessen die globale Wärmespur divergieren.
Der Text entfernt diese Fragen nicht durch Vergrößerung.

\subsection{2. Der kleinste Gegencheck zur
Primzahl-Automatik}\label{der-kleinste-gegencheck-zur-primzahl-automatik}

Nehmen wir einen gerichteten Knoten mit zwei erlaubten Schleifen a und
b. Neben a und b ist auch ab ein primitiver periodischer Weg: Es ist
keine Potenz eines kürzeren Wortes. Ebenso entstehen weitere gemischte
primitive Wörter. Setzen wir nachträglich L(a)=log 2, L(b)=log 3, so hat
ab die Länge log 6. Die 6 ist keine Primzahl. „Primitiver Weg`` bedeutet
nicht „arithmetische Primzahl``.

Schon formal unterscheiden sich die zugehörigen erzeugenden Funktionen:

\[\adjustbox{max width=\linewidth}{$\displaystyle 
Z_{\rm Wege}(x,y)=\frac1{1-x-y},\qquad
Z_{\rm zwei\ Primarten}(x,y)=\frac1{(1-x)(1-y)}.
$}\]

Der Koeffizient von xy ist links 2 und rechts 1. Links werden die beiden
Wörter ab und ba gezählt; rechts eine kommutative Besetzung. Im
periodischen Eulerprodukt bilden ab und ba eine zyklische Klasse, aber
diese ist ein \textbf{neuer primitiver Faktor}. Die Abweichung
verschwindet damit nicht.

Eine kommutative freie Halbgruppe über vorgegebenen Primarten
reproduziert eindeutige Faktorzerlegung. Sie leitet aber weder die
natürliche arithmetische Markierung noch die logarithmischen Längen oder
die passende Spur her. Eine gekoppeltere Geometrie müsste gemischte
Bahnen mit den richtigen Amplituden, Relationen oder nachgewiesenen
Auslöschungen behandeln. Dies ist kein allgemeiner Ausschluss von
Schleifenmodellen mit Interferenz.

Der Unterschied ist in den Originalarbeiten sichtbar: Kuipers, Hummel
und Richter konstruieren Quantengraphen mit dem passenden oszillierenden
Anteil der Nullstellendichte, weisen aber auf den anderen glatten Anteil
und damit das andere Spektrum hin. Das ist ein nützlicher Vorläufer,
kein RH-Operator. \href{https://arxiv.org/abs/1307.6055}{Originalarbeit,
Phys. Rev.~Lett. 112, 070406}

Graphische Eulerprodukte besitzen ihre eigene Theorie. Auch eine aus
einer Graph-Zetafunktion bestimmbare Größe ist nicht automatisch
effizient auslesbar. Storm trennt in seiner Arbeit ausdrücklich
berechenbare Zeta-Darstellung und teure Auswertung bestimmter
Graphinformationen.
\href{https://arxiv.org/abs/0708.1923}{Originalarbeit zu
Edge-Zetafunktionen}

